% VI/39 reader bridge -- additive mathematical exposition.
\paragraph{What a Weil divisor measures.}
Let $X$ be a Noetherian, normal, integral scheme with function field
$K(X)$.  A \emph{prime Weil divisor} is an integral closed subscheme
of codimension one.  A Weil divisor is a finite formal integer sum
\[
 D=\sum_Z n_Z[Z],\qquad n_Z\in\mathbb Z,
\]
where $Z$ runs over prime divisors.  ``Formal'' matters: the coefficients
record orders of vanishing or poles, not just a union of closed sets.
The divisor is \emph{effective} when all $n_Z\geq0$; its support
consists of the $Z$ having nonzero coefficients.  Only finitely many
coefficients may be nonzero.  Unlike a divisor of a single global
function, an arbitrary Weil divisor need not have one defining equation.

\paragraph{Why normality enters the definition of principal divisors.}
For a codimension-one point $\eta_Z$, normality and Noetherianness
imply that the one-dimensional local domain
$\mathcal O_{X,\eta_Z}$ is a discrete valuation ring.  It has a
valuation
$\operatorname{ord}_Z\colon K(X)^\times\to\mathbb Z$.
Consequently every nonzero rational function $f$ defines
\[
 \operatorname{div}(f)=\sum_Z\operatorname{ord}_Z(f)[Z].
\]
The sum is finite.  If $f$ is a unit near the generic point of $Z$,
its coefficient is zero; if a local uniformizer occurs twice in $f$,
its coefficient is $2$.  Multiplication of rational functions becomes
addition of their principal divisors:
$\operatorname{div}(fg)=\operatorname{div}(f)+\operatorname{div}(g)$.
Without suitable normality hypotheses, local rings at codimension-one
points need not be discrete valuation rings, so this valuation formula
cannot simply be used unchanged.

\paragraph{Worked example 1: three lines and a rational function.}
Take $X=\mathbb A_k^2=\operatorname{Spec}k[x,y]$ and
$f=x^2y/(x-y)\in k(x,y)^\times$.  The irreducible equations
$x$, $y$, and $x-y$ define three different codimension-one prime
divisors.  At the generic point of $V(x)$, both $y$ and $x-y$
are units, so the order is two.  At $V(y)$, the order is one.
At $V(x-y)$, the denominator has order one, producing a pole.
No other irreducible polynomial enters the factorization.  Therefore
\[
 \operatorname{div}\!\left(\frac{x^2y}{x-y}\right)
    =2[V(x)]+[V(y)]-[V(x-y)].
\]
A nonzero rational function can thus have a divisor with both
positive and negative coefficients.  In particular, the divisor
remembers a double zero along one entire line, not merely the points
where that line meets the other two.

\paragraph{Worked example 2: a nonconstant unit has no divisor.}
For the affine hyperbola
$X=\operatorname{Spec}k[x,y]/(xy-1)$, the coordinate $x$ is
invertible with inverse $y$.  Hence it has no zeros and no poles
at any codimension-one point of $X$:
\[
 \operatorname{div}(x)=0\quad\text{on }X.
\]
This does \emph{not} say that $x$ is constant.  It says only that
$x$ belongs to the group $\Gamma(X,\mathcal O_X)^\times$ of global
units.  Removing the points at which a function has a zero or pole
may change its divisor, because divisors depend on the scheme on
which the rational function is considered.

\paragraph{Worked example 3: zero and pole on the projective line.}
Write $\mathbb P_k^1=\operatorname{Proj}k[X,Y]$ and use the
rational coordinate $t=X/Y$ on the chart $Y\ne0$.  It vanishes
once at $P_0=[0:1]$.  On the other chart put $s=Y/X$;
then $t=1/s$, giving a simple pole at $P_\infty=[1:0]$.
Hence
\[
 \operatorname{div}(t)=[P_0]-[P_\infty].
\]
This calculation anticipates the degree-zero property of principal
divisors on a complete nonsingular curve.  It also makes clear why
the choice of an affine chart can hide the pole at infinity.

\paragraph{From equations to divisors.}
In $\mathbb P_k^2$, let $F$ and $G$ be nonzero homogeneous
polynomials of the same degree.  The quotient $F/G$ is a rational
function on projective space: multiplying all homogeneous coordinates
by $\lambda$ multiplies numerator and denominator by the same
$\lambda^d$.  Its divisor is the difference of the corresponding
hypersurface divisors, with all irreducible factors counted with
multiplicities.  For example,
\[
 \operatorname{div}(X/Y)=[V_+(X)]-[V_+(Y)].
\]
The bracketed terms denote codimension-one subschemes, whereas
the ratio is a rational function; mixing these two roles is a frequent
source of incorrect projective arguments.
